% REPLACES: main.tex lines 1815-1880 ("\subsubsection{Logistic noise}" to "posterior mass sits relative to the common bounds \(m,M\).")
% RETAINS LABELS: logistic-noise, result-proposition-a-5, eq:paper-a-21,
% eq:paper-a-22, eq:paper-a-23. No label is removed.
% NEW SYMBOLS: mathcal C_L, L, mu_f(x;pi), bar mu_f(pi), mathcal E_f, mathsf E_f.
% The symbols denote the density class, score bound, full-order posterior,
% posterior supremum, favorable flow set, and entry. Existing alpha_j refers to this set.
% DEPENDS ON: the quality-class Proposition 2, prior pi, zero preparation cost,
% kappa(pi), and the general nonemptiness argument in the main-model revision.
% OWNERSHIP: use the exact condition below; the Laplace-specific sufficient
% condition based on alpha_H > 1/2 does not extend to every class member.
% NUMBERS: none. CITES: ref-Blackwell1953 already exists.
% SCOPE: the common-variance comparison below is a numerical diagnostic.

\subsubsection{Bounded-score noise}\label{logistic-noise}

Let \(\mathcal C_L\) contain the positive, absolutely continuous densities
\(f\) with \(|f'|\le Lf\) almost everywhere, where \(0<L<1\).
Equivalently, \(\log f\) is Lipschitz with constant \(L\).
Set \(b=1/L\).
The bounds \(m(\pi),M(\pi)\) and
\(\kappa(\pi)=\min\{m(\pi),1-M(\pi)\}\) use this value of \(b\).
For full orders, define

\[
\mu_f(x;\pi)=
\frac{\pi f(x-1)}{\pi f(x-1)+(1-\pi)f(x+1)},\qquad
\bar\mu_f(\pi)=\operatorname*{ess\,sup}_{x\in\mathbb R}\mu_f(x;\pi).
\]

The supremum satisfies \(\pi<\bar\mu_f(\pi)\le M(\pi)\).
It need not be attained at a finite flow.
It can also be strictly below the uniform bound \(M(\pi)\).

\Needspace{8\baselineskip}

\protect\phantomsection\label{result-proposition-a-5}{\textbf{Proposition
A.5 (bounded-score noise; analytical).}} \emph{Replace Laplace noise by
any \(f\in\mathcal C_L\).
Keep the quality-class assumptions of
\hyperref[result-proposition-2]{Proposition 2}, with preparation costs
\(C\in\{0,c_H\}\).
Replace the upper posterior in \eqref{eq:paper-a2} by \(\bar\mu_f(\pi)\), so that}

\[
B_{r_0}(\pi)<c_H<B_{r_1}(\bar\mu_f(\pi)).
\]

\emph{Retain \eqref{eq:paper-a3} with \(b=1/L\), namely}

\[
\Delta_T(r_0)<k<(1-L)\rho\kappa(\pi)\Delta_T(r_1).
\]

\emph{Then the weak economy has unique zero orders and entry \(\rho\).
The strong economy has unique full orders \((1,-1)\) and an informative price.
Let}

\[
\begin{aligned}
\tau&=\frac{c_H-g_L(r_1)}{g_H(r_1)-g_L(r_1)},
&\mathcal E_f&=\{x:\mu_f(x;\pi)\ge\tau\},\\
\alpha_H&=\int_{\mathcal E_f}f(x-1)\,dx,
&\alpha_L&=\int_{\mathcal E_f}f(x+1)\,dx.
\end{aligned}
\]

\emph{Total entry satisfies}

\[
\mathsf E_f=\rho+(1-\rho)[\pi\alpha_H+(1-\pi)\alpha_L]>\rho.
\]

\emph{High-class ownership rises exactly when
\([\rho+(1-\rho)\alpha_H]W_H(r_1)>\rho W_H(r_0)\).
The condition holds automatically in the point-mass case with \(W_H=1\).
At a later strength, part (iii) retains the uniform bound
\(B_{r_2}(M(\pi))<c_H\) and its stated order condition.
These comparisons hold on a nonempty open set of remaining primitives for fixed \(f\) and zero preparation cost.}

\emph{Proof.}
The score bound implies
\(|\log f(y)-\log f(z)|\le L|y-z|\).
Hence \(f(y)\ge f(z)e^{-L|y-z|}\), so \(f(z)\le L/2\) after integration over \(y\).
Also \(\int|f'|\le L\), and therefore \(f\in W^{1,1}\).
The density is positive, bounded, and continuous.
The null-set arguments in Online Appendix A.1 thus apply under every mixed order and every unilateral deviation.
\par
For any feasible \(q,q'\), the density ratio lies in \([e^{-2L},e^{2L}]\).
Integration over the order distributions and Bayes' rule give
\(m(\pi)\le\mu_X\le M(\pi)\).
Conditional expectation gives the same bounds for the price posterior.

The zero preparation cost gives \(e\ge\rho\) at every belief.
The price inversion and the residual calculation in Online Appendix A.3 therefore apply without change.
They use only \(t_0,t_j,g_j,\Delta_T\), independence, and the positive entry floor.
In particular,

\[
\rho\kappa(\pi)\Delta_T\le A_H,A_L\le\Delta_T.
\]

The translation identity, Fubini argument, and product rule in Online Appendix A.3 already allow this general \(f\).
They give \(|F_j'|\le LF_j\) and

\[
U_j'(s)\ge(1-sL)F_j(s)-k
\ge(1-L)\rho\kappa(\pi)\Delta_T(r_1)-k>0.
\]

These steps differentiate the density; they do not differentiate entry.
At \(r_0\), every correctly signed nonzero order earns at most
\(s[\Delta_T(r_0)-k]<0\).
A wrong sign also loses.
Thus zero is the unique best response, including among mixed orders.
The posterior equals \(\pi\), and the lower cost restriction excludes expensive preparation.
The constant competitive price constructs the weak outcome.
\par
At \(r_1\), the positive derivative makes the unit magnitude strictly optimal for each class.
It also excludes every mixture with another order.
Define entry and price by \eqref{eq:paper-9}, with posterior \(\mu_f(x;\pi)\).
The price is strictly increasing in the posterior, so its measurable inverse proves price sufficiency.
This step needs no monotonicity of the posterior in flow.
The construction gives competitive prices and optimal preparation, and the global bound verifies every order deviation.

Gross preparation profit falls with incumbent strength at every fixed belief.
The cost window therefore gives \(\pi<\tau<\bar\mu_f(\pi)\).
The set \(\{\mu_f>\tau\}\) has positive Lebesgue measure by the definition of the essential supremum.
Positivity of both conditional densities gives \(\alpha_H,\alpha_L>0\).
On \(\mathcal E_f\), Bayes' rule gives

\[
f(x-1)\ge
\frac{(1-\pi)\tau}{\pi(1-\tau)}f(x+1),\qquad
\alpha_H\ge\frac{(1-\pi)\tau}{\pi(1-\tau)}\alpha_L>\alpha_L.
\]

The entry formula and the stated ownership condition follow.
The likelihood ratio cannot be constant: integration would make that constant one, and then \(f\) would be two-periodic.
A positive integrable density cannot be two-periodic.
The strong price experiment is therefore informative and strictly Blackwell-dominates the weak constant experiment.
Since both conditional densities integrate to one, their nonconstant ratio exceeds one on a set of positive measure.
Thus \(\bar\mu_f(\pi)>\pi\).
\par
For part (iii), the order bound again forces full orders, while
\(\mu_X\le M(\pi)\) excludes expensive preparation.

The nonemptiness construction below uses a positive class profit gap where the target spread tends to zero.
Replace its upper posterior by \(\bar\mu_f(\pi)>\pi\).
The cost window remains nonempty, and \(\kappa(\pi)>0\) leaves a nonempty order-cost window.
\par
For ownership, first fix an interior threshold at a continuity point of the posterior distribution.
Its extra entry probability stays positive near that threshold.
Choose the two incumbent strengths sufficiently near each other that their high-class win probabilities satisfy the stated ownership condition.
Continuity preserves every strict inequality on an open neighborhood with the zero preparation cost fixed.
This is the construction of
\hyperref[nonemptiness-for-every-value-gap]{the nonemptiness argument} with the new upper posterior.
\(\square\)

\emph{Logistic example.}
The density

\[
f_{\mathrm{log}}(z)=\frac{1}{4b}\operatorname{sech}^{2}\left(\frac z{2b}\right)
=\frac{e^{-z/b}}{b(1+e^{-z/b})^2},\qquad b>1,
\tag{A.21}\label{eq:paper-a-21}
\]

has log derivative \(-\tanh(z/2b)/b\), so it belongs to \(\mathcal C_{1/b}\).
Under full orders its posterior obeys

\[
\operatorname{logit}\mu_X
=\operatorname{logit}\pi
+2\log\cosh\left(\frac{x+1}{2b}\right)
-2\log\cosh\left(\frac{x-1}{2b}\right).
\tag{A.22}\label{eq:paper-a-22}
\]

The derivative is
\([\tanh((x+1)/2b)-\tanh((x-1)/2b)]/b>0\).
The posterior spans \((m(\pi),M(\pi))\), with neither endpoint attained.
Thus \(\bar\mu_f(\pi)=M(\pi)\), and the original upper restriction applies.
\par
Set \(A=e^{1/b}\) and
\(w=\sqrt{\tau(1-\pi)/[(1-\tau)\pi]}\).
With \(y=e^{x/b}\), the square root of the likelihood ratio is
\((Ay+1)/(y+A)\).
The equation \(w=(Ay+1)/(y+A)\) gives

\[
\begin{aligned}
x^*_{\mathrm{log}}&=b\log\frac{Aw-1}{A-w},\\
\alpha_H^{\mathrm{log}}&=\frac1{1+e^{(x^*_{\mathrm{log}}-1)/b}},\qquad
\alpha_L^{\mathrm{log}}=\frac1{1+e^{(x^*_{\mathrm{log}}+1)/b}}.
\end{aligned}
\tag{A.23}\label{eq:paper-a-23}
\]

These formulas hold for \(m(\pi)<\tau<M(\pi)\).
At \(\pi=1/2\), they reduce to the benchmark formulas.
Other members include hyperbolic-secant noise, suitably scaled Student noise, and Gaussian noise plus an independent Laplace component.

Laplace and logistic noise are compared at a common scale parameter \(b\).
They do not have a common variance.
At a common scale, the logistic law equals the Laplace law plus independent noise.
Thus the Laplace flow experiment Blackwell-dominates the logistic experiment at every order profile, pure or mixed
(\citeproc{ref-Blackwell1953}{Blackwell 1953}).
Price sufficiency transfers this comparison to prices, including the full-order equilibria compared here.
Online Appendix A.6 proves the decomposition.
\par
At a common variance, neither experiment dominates in the numerical comparison.
This last comparison is a numerical diagnostic.
Dominance does not order entry: \(\Pr(\mu_X\ge\tau)\) is not convex in the posterior.
The entry comparison concerns where posterior mass sits relative to the common bounds \(m(\pi),M(\pi)\).
